\section{Comparison Methods and Baselines}\label{sec:comparators}
\subsection{Protected references and component removals}
The final roster compares base ARX, linear correction and RBF correction under matched inputs. Point rules directly threshold $a_t$. Return rules retain alarm using a one-point residual condition, the 41-point window, or the 5/41 dual window, generally with ten-point dwell. The no-dwell variant requires one qualifying point. Base long versus base dual and RBF long versus RBF dual isolate the short window; RBF dual versus its no-dwell variant isolates persistence. Base dual versus RBF dual is the primary correction comparison. Every variant keeps its declared entry score, threshold and support.

The RBF-following ablation applies the correction trained on free-running features to measured-history predictions without refitting. It tests that disclosed assembly, not an optimally fitted following model. Table~\ref{tab:audit} retains all 17 methods.

\subsection{Ordinary prediction and accumulation controls}
Persistence predicts the last measured state. The causal rolling mean predicts from the preceding 21 observations, excluding the current target. A one-sided CUSUM uses $G_t=\max(0,G_{t-1}+a_t-k)$, with $k$ equal to the training-score mean plus its standard deviation. Its calibrated statistic has no reset on alarm in this registered implementation. Its persistence is therefore a deliberate comparator behavior, not a claim that all CUSUM schemes cannot recover. All controls use identical eligible episode support and cannot access labels or truth.

The IMU controls are training mean, persistence and AR(4), alongside the CNNs specified in Section~\ref{sec:experiments}. AR fits standardized channels with ridge $10^{-3}$ on mean normal equations, leaving its intercept unpenalized. Deterministic controls run once and neural models use all three seeds.

\subsection{Recent published methods and correspondence}
The recent-paper comparison window is 16 October 2023--16 October 2026, with preference for 2024--2026. CASPER (2024) and Edge CNN (2025) are independent architectural implementations; ModernTCN, iTransformer and TimeMixer (ICLR 2024) execute author model code; MVT-Flow (IEEE T-RO 2024) remains unexecuted. Publication recency does not establish comparability or a current state-of-the-art ranking. Classical AR, persistence, rolling means and CUSUM remain explanatory controls. Older sources support observer, uncertainty and input-dependent-threshold foundations; they are not recent detector baselines.



The CNNs use independent implementations; author forecasters use unchanged model sources with our causal adapter, preprocessing and budget. These normal-development measurements are not pooled into the protected-command attack ranking. MVT-Flow remains an unexecuted offline candidate; its symmetric convolutions, whole-execution objective and source correspondence require a separate task contract. No complete published attack-detector reproduction is claimed.
